\documentclass[11pt,openany]{book}
\usepackage[a4paper,margin=2.45cm]{geometry}
\usepackage{amsmath,amssymb,amsfonts,bm,mathtools}
\usepackage{booktabs}
\usepackage{enumitem}
\usepackage{hyperref}
\usepackage{xcolor}
\usepackage{microtype}
\hypersetup{colorlinks=true,linkcolor=blue,urlcolor=blue,citecolor=blue,pdftitle={From Particles to Computational Grids},pdfauthor={}}
\setcounter{tocdepth}{-1}
\setlength{\parindent}{2em}
\setlength{\parskip}{0.35em}
\newcommand{\dd}{\,\mathrm d}
\newcommand{\ee}{\mathrm e}
\newcommand{\ii}{\mathrm i}
\newcommand{\R}{\mathbb R}
\newcommand{\Order}{\mathcal O}
\newcommand{\Lop}{\mathcal L}
\newcommand{\Nop}{\mathcal N}
\newcommand{\Ksp}{\mathcal K}
\newcommand{\avg}[1]{\left\langle #1\right\rangle}
\title{From Particles to Computational Grids\\\large An Introduction to Computational Mechanics for Readers with a Mathematical Background}
\author{}
\date{}
\begin{document}
\maketitle
\tableofcontents

\chapter*{Preface}
Computational mechanics is not merely the act of placing physical equations inside a computer. It is a continuous chain of information reduction:
\[
\text{many-body dynamics}\to\text{statistical description}\to\text{continuum model}\to\text{discrete equations}\to\text{matrix system}\to\text{engineering computation}.
\]
The physical system, mathematical model, discrete model, and algebraic system are distinct objects. Physical modeling determines what is neglected; numerical discretization determines what is approximated; linear algebra determines how the discrete problem is solved; and the grid connects geometry, boundary conditions, and local scales to the discrete system.

\part{Theoretical Framework}

\chapter*{From Many-Particle Dynamics to Kinetic Equations}
Consider a system of $N$ particles. The position and momentum of particle $i$ are
\[
\bm x_i,\qquad \bm p_i.
\]
The complete state is a point in a $6N$-dimensional phase space:
\[
\Gamma=(\bm x_1,\dots,\bm x_N,\bm p_1,\dots,\bm p_N).
\]
If the Hamiltonian is $H$,
\[
\dot{\bm x}_i=\frac{\partial H}{\partial \bm p_i},\qquad
\dot{\bm p}_i=-\frac{\partial H}{\partial \bm x_i}.
\]
For macroscopic matter, direct particle-by-particle integration is infeasible. Introduce the phase-space density $\rho_N(\Gamma,t)$. Conservation of probability gives
\[
\frac{\partial \rho_N}{\partial t}+\nabla_\Gamma\cdot(\rho_N\dot\Gamma)=0.
\]
For Hamiltonian flow, $\nabla_\Gamma\cdot\dot\Gamma=0$, so
\[
\frac{\partial \rho_N}{\partial t}+\dot\Gamma\cdot\nabla_\Gamma\rho_N=0,
\]
or
\[
\frac{\partial \rho_N}{\partial t}+\{\rho_N,H\}=0.
\]
This is the Liouville equation. It preserves phase-space volume but remains $6N$-dimensional.

Define the one-particle distribution function $f(\bm x,\bm v,t)$. Integrating out the remaining particles produces the BBGKY hierarchy
\[
f_1\to f_2\to f_3\to\cdots,
\]
so closure assumptions are unavoidable.

A general kinetic equation is
\[
\frac{\partial f}{\partial t}+\bm v\cdot\nabla_{\bm x}f+\frac{\bm F}{m}\cdot\nabla_{\bm v}f=C[f].
\]
For collisionless charged particles,
\[
\frac{\partial f_s}{\partial t}+\bm v\cdot\nabla_{\bm x}f_s+\frac{q_s}{m_s}(\bm E+\bm v\times\bm B)\cdot\nabla_{\bm v}f_s=0.
\]
The charge and current densities are
\[
\rho_c=\sum_s q_s\int f_s\dd\bm v,\qquad
\bm J=\sum_s q_s\int \bm v f_s\dd\bm v.
\]
For dilute neutral gases, the Boltzmann equation is
\[
\frac{\partial f}{\partial t}+\bm v\cdot\nabla_{\bm x}f+\frac{\bm F}{m}\cdot\nabla_{\bm v}f=Q(f,f),
\]
with
\[
Q(f,f)=\int B(f'f_*'-ff_*)\dd\Omega\dd\bm v_*.
\]
The molecular-chaos closure is
\[
f_2(\bm v,\bm v_*)\approx f(\bm v)f(\bm v_*).
\]
The BGK model replaces the collision integral by
\[
C[f]=\frac{f^{\mathrm{eq}}-f}{\tau},
\]
where
\[
f^{\mathrm{eq}}=n\left(\frac{m}{2\pi k_BT}\right)^{3/2}\exp\left[-\frac{m|\bm v-\bm u|^2}{2k_BT}\right].
\]
For long-range Coulomb interactions, Landau or Fokker--Planck operators describe velocity-space drift and diffusion.

\chapter*{From Velocity Moments to Continuum Equations}
The number density, mass density, and bulk velocity are
\[
n=\int f\dd\bm v,\qquad \rho=mn,\qquad \rho\bm u=m\int \bm v f\dd\bm v.
\]
With $\bm c=\bm v-\bm u$,
\[
\bm P=m\int \bm c\bm c f\dd\bm v,\qquad
\bm q=\frac{m}{2}\int |\bm c|^2\bm c f\dd\bm v.
\]
The zeroth moment gives
\[
\boxed{\frac{\partial\rho}{\partial t}+\nabla\cdot(\rho\bm u)=0}.
\]
The first moment gives
\[
\frac{\partial(\rho\bm u)}{\partial t}+\nabla\cdot(\rho\bm u\bm u+\bm P)=\rho\bm b.
\]
Writing $\bm P=p\bm I-\bm\tau$,
\[
\frac{\partial(\rho\bm u)}{\partial t}+\nabla\cdot(\rho\bm u\bm u)=-\nabla p+\nabla\cdot\bm\tau+\rho\bm b.
\]
For a Newtonian fluid,
\[
\bm\tau=\mu[\nabla\bm u+(\nabla\bm u)^T]+\lambda(\nabla\cdot\bm u)\bm I.
\]
For an incompressible fluid,
\[
\nabla\cdot\bm u=0,
\]
\[
\rho\left(\frac{\partial\bm u}{\partial t}+\bm u\cdot\nabla\bm u\right)=-\nabla p+\mu\nabla^2\bm u+\rho\bm b.
\]
The energy equation is
\[
\frac{\partial(\rho E)}{\partial t}+\nabla\cdot[(\rho E+p)\bm u]=\nabla\cdot(\bm\tau\cdot\bm u)-\nabla\cdot\bm q+\rho\bm b\cdot\bm u,
\]
where $E=e+|\bm u|^2/2$ and $\bm q=-k\nabla T$.

\chapter*{Maxwell's Equations, Linearization, and Perturbation Operators}
Maxwell's equations are
\[
\nabla\cdot\bm E=\frac{\rho_c}{\varepsilon_0},\qquad \nabla\cdot\bm B=0,
\]
\[
\frac{\partial\bm B}{\partial t}+\nabla\times\bm E=0,
\]
\[
\frac{\partial\bm E}{\partial t}-c^2\nabla\times\bm B=-\frac{\bm J}{\varepsilon_0}.
\]
They imply charge conservation:
\[
\frac{\partial\rho_c}{\partial t}+\nabla\cdot\bm J=0.
\]
The Lorentz force density is
\[
\bm f_L=\rho_c\bm E+\bm J\times\bm B.
\]

Write a nonlinear coupled model as $\Nop(\bm U)=0$. Near a base state,
\[
\bm U=\bm U_0+\epsilon\bm u,\qquad 0<\epsilon\ll1.
\]
Then
\[
\Nop(\bm U_0+\epsilon\bm u)=\Nop(\bm U_0)+\epsilon D\Nop[\bm U_0](\bm u)+\frac{\epsilon^2}{2}D^2\Nop[\bm U_0](\bm u,\bm u)+\cdots.
\]
If $\Nop(\bm U_0)=0$,
\[
\Lop\bm u=0,\qquad \Lop:=D\Nop[\bm U_0].
\]
For nonlinear constitutive laws $\bm D=\bm D(\bm E)$ and $\bm B=\bm B(\bm H)$,
\[
\bm D(\bm E_0+\epsilon\bm e)=\bm D(\bm E_0)+\epsilon\left.\frac{\partial\bm D}{\partial\bm E}\right|_{\bm E_0}\bm e+\Order(\epsilon^2),
\]
\[
\bm B(\bm H_0+\epsilon\bm h)=\bm B(\bm H_0)+\epsilon\left.\frac{\partial\bm B}{\partial\bm H}\right|_{\bm H_0}\bm h+\Order(\epsilon^2).
\]
For Vlasov--Maxwell, let
\[
f_s=f_{s0}+\epsilon f_{s1},\qquad \bm E=\epsilon\bm E_1,\qquad \bm B=\bm B_0+\epsilon\bm B_1.
\]
The linearized equation is
\[
\frac{\partial f_{s1}}{\partial t}+\bm v\cdot\nabla_{\bm x}f_{s1}+\frac{q_s}{m_s}(\bm v\times\bm B_0)\cdot\nabla_{\bm v}f_{s1}
=-\frac{q_s}{m_s}(\bm E_1+\bm v\times\bm B_1)\cdot\nabla_{\bm v}f_{s0}.
\]
With a plane-wave ansatz,
\[
\bm u_1=\widehat{\bm u}\exp[\ii(\bm k\cdot\bm x-\omega t)],
\]
we have $\partial_t\mapsto-\ii\omega$ and $\nabla\mapsto\ii\bm k$. The dispersion relation follows from
\[
\det\bm M(\omega,\bm k)=0.
\]

\chapter*{Scale Analysis and Asymptotic Expansions}
Typical nondimensional parameters are
\[
\mathrm{Kn}=\frac{\lambda}{L},\qquad \mathrm{Re}=\frac{\rho UL}{\mu},\qquad \mathrm{Ma}=\frac{U}{c}.
\]
A Taylor expansion is
\[
f(x+h)=f(x)+hf'(x)+\frac{h^2}{2}f''(x)+\cdots.
\]
An asymptotic expansion is more general. If $\phi_{n+1}(\epsilon)=o(\phi_n(\epsilon))$ and
\[
f(\epsilon)-\sum_{n=0}^{N}a_n\phi_n(\epsilon)=o(\phi_N(\epsilon)),
\]
then
\[
f(\epsilon)\sim\sum_{n=0}^{\infty}a_n\phi_n(\epsilon).
\]
The infinite series need not converge. Scales may include
\[
1,\quad \epsilon^{1/2},\quad \epsilon\log\epsilon,\quad \ee^{-1/\epsilon}.
\]
If terms first decrease and then increase, the best approximation is obtained by optimal truncation near the smallest term.

A regular perturbation preserves differential order and boundary-condition structure. A singular perturbation occurs when the limit lowers the equation order, for example
\[
\epsilon u''+u'=0.
\]

\section*{Plasma Example: Small Debye Length and the Quasineutral Limit}
Consider
\[
-\epsilon^2\frac{\partial^2\phi}{\partial x^2}=n_i-n_e,\qquad \epsilon=\frac{\lambda_D}{L}\ll1.
\]
Formally setting $\epsilon=0$ gives quasineutrality:
\[
n_i=n_e.
\]
If the wall potential is incompatible with the bulk solution, the problem is singular. Write
\[
\phi^{\mathrm{out}}=\phi_0+\epsilon\phi_1+\epsilon^2\phi_2+\cdots.
\]
Near the wall, introduce $\xi=x/\epsilon$, so
\[
\frac{\partial}{\partial x}=\frac{1}{\epsilon}\frac{\partial}{\partial\xi}.
\]
The inner equation becomes
\[
-\frac{\partial^2\phi^{\mathrm{in}}}{\partial\xi^2}=n_i-n_e.
\]
The Debye sheath has thickness $\Order(\lambda_D)$. It is thin macroscopically but essential for enforcing the wall potential and accommodating charge separation.

\part{Numerical Discretization}

\chapter*{Consistency, Stability, and Modified Equations}
Consider $u_t+au_x=0$. The forward and centered differences satisfy
\[
\frac{u_i^{n+1}-u_i^n}{\Delta t}=u_t+\Order(\Delta t),
\]
\[
\frac{u_{i+1}-u_{i-1}}{2\Delta x}=u_x+\Order(\Delta x^2).
\]
For a Fourier error mode $e_i^n=\widehat e^n\ee^{\ii kx_i}$,
\[
\widehat e^{n+1}=G(k)\widehat e^n.
\]
Stability requires $|G(k)|\le1$. With $\nu=a\Delta t/\Delta x$, first-order upwinding gives
\[
u_i^{n+1}=(1-\nu)u_i^n+\nu u_{i-1}^n.
\]
Its modified equation is
\[
u_t+au_x=\frac{a\Delta x}{2}(1-\nu)u_{xx}+\cdots.
\]

\chapter*{Godunov's Theorem and High-Resolution Schemes}
Godunov's theorem states that a linear, monotone, consistent scheme for scalar linear advection can be at most first-order accurate. High-order linear reconstruction generally requires weights of alternating sign, whereas monotonicity requires a convex-combination structure.

Modern schemes retain conservation, introduce nonlinear adaptation, recover high order in smooth regions, and reduce toward first order near discontinuities. A finite-volume scheme has the form
\[
\frac{\dd \bar U_i}{\dd t}=-\frac{1}{\Delta x}(\widehat F_{i+1/2}-\widehat F_{i-1/2}).
\]
MUSCL reconstruction may use
\[
U_{i+1/2}^{L}=U_i+\frac12\sigma_i,
\]
with nonlinear limiters such as minmod, van Leer, or Superbee. TVD schemes require
\[
TV(U^{n+1})\le TV(U^n),\qquad TV(U)=\sum_i|U_{i+1}-U_i|.
\]
ENO and WENO methods select or weight candidate stencils so that smooth regions receive high-order reconstruction while discontinuous stencils are suppressed.

\chapter*{Dispersion Relations and Wave-Propagation Error}
For $u=\widehat u\ee^{\ii(kx-\omega t)}$, exact advection gives
\[
\omega=ak,
\]
with
\[
c_p=\frac{\omega}{k}=a,\qquad c_g=\frac{\dd\omega}{\dd k}=a.
\]
A discrete scheme produces $\omega_{\mathrm{num}}(k)$. Errors in its real part cause numerical dispersion; a negative imaginary part causes numerical dissipation.

\chapter*{Finite Differences, Finite Volumes, Finite Elements, and Spectral Methods}
Finite differences approximate derivatives directly. Finite volumes discretize
\[
\frac{\dd}{\dd t}\int_{\Omega_i}U\dd V+\int_{\partial\Omega_i}\bm F\cdot\bm n\dd S=\int_{\Omega_i}S\dd V.
\]
Finite elements use weak formulations. For $-\nabla^2u=f$,
\[
\int_\Omega\nabla v\cdot\nabla u\dd V=\int_\Omega vf\dd V+\text{boundary terms}.
\]
Spectral methods use global basis functions and are highly accurate for smooth solutions.

\part{Matrix Solvers}

\chapter*{From Discrete Operators to Algebraic Systems}
Discretization produces
\[
\bm A\bm x=\bm b
\]
or
\[
\bm R(\bm U)=0.
\]
Newton's method uses
\[
\bm J(\bm U^k)\Delta\bm U=-\bm R(\bm U^k),\qquad \bm U^{k+1}=\bm U^k+\Delta\bm U.
\]

\chapter*{LU and LU-SGS}
LU factorization is
\[
\bm A=\bm L\bm U,
\]
followed by
\[
\bm L\bm y=\bm b,\qquad \bm U\bm x=\bm y.
\]
Dense LU costs $O(N^3)$ operations and $O(N^2)$ storage. LU-SGS approximates
\[
\bm A\approx(\bm D+\bm L)\bm D^{-1}(\bm D+\bm U).
\]

\chapter*{Krylov Methods, Arnoldi, GMRES, and BiCGSTAB}
The Krylov subspace is
\[
\Ksp_m(\bm A,\bm r_0)=\operatorname{span}\{\bm r_0,\bm A\bm r_0,\dots,\bm A^{m-1}\bm r_0\}.
\]
Arnoldi gives
\[
\bm A\bm V_m=\bm V_{m+1}\overline{\bm H}_m.
\]
GMRES solves
\[
\bm x_m=\arg\min_{\bm x\in\bm x_0+\Ksp_m}\|\bm b-\bm A\bm x\|_2.
\]
BiCGSTAB uses short recurrences and fixed memory, although its residual history may oscillate.

\chapter*{Preconditioning and Algebraic Multigrid}
The condition number is
\[
\kappa(\bm A)=\|\bm A\|\,\|\bm A^{-1}\|.
\]
A preconditioner aims to make $\bm M^{-1}\bm A$ easier to solve. AMG constructs coarse levels from matrix connectivity and applies smoothing, restriction, coarse-grid correction, prolongation, and postsmoothing.

\chapter*{Method of Moments and Dense Integral Equations}
The method of moments begins with
\[
\Lop J=g.
\]
With
\[
J=\sum_{n=1}^N I_nf_n,
\]
testing gives
\[
\avg{w_m,\Lop J}=\avg{w_m,g},
\]
and hence
\[
\bm Z\bm I=\bm V.
\]
Because the interaction kernel is long-ranged, $\bm Z$ is usually dense.

\part{Grid Techniques}

\chapter*{The Grid as the Interface Between Geometry and Algebra}
The grid controls geometric representation error, discretization error, stability restrictions, matrix sparsity, and parallel communication. Body-fitted grids conform to boundaries and resolve wall layers well, but become difficult under complex geometry or large motion.

\chapter*{Immersed Boundary Methods}
Immersed boundary methods solve on a simple background grid and impose the body through
\[
\rho\left(\frac{\partial\bm u}{\partial t}+\bm u\cdot\nabla\bm u\right)=-\nabla p+\mu\nabla^2\bm u+\bm f_{\mathrm{IB}}.
\]
They simplify meshing and moving-geometry treatment, but wall accuracy and strict conservation require care.

\chapter*{Quadtrees, Octrees, and Adaptive Refinement}
Quadtrees and octrees refine cells recursively. Engineering constraints include 2:1 balance, coarse-fine flux treatment, conservative data transfer, and dynamic load balancing.

\chapter*{Moving Meshes and the Geometric Conservation Law}
With mesh velocity $\bm w$,
\[
\frac{\dd}{\dd t}\int_{\Omega(t)}U\dd V+\int_{\partial\Omega(t)}[\bm F(U)-U\bm w]\cdot\bm n\dd S=0.
\]
The geometric conservation law requires mesh motion alone not to change a constant physical state.

\chapter*{Overset Grids}
Overset grids superimpose independent mesh blocks. Donor cells provide interpolated data to receptor cells. The main challenges are hole cutting, donor search, conservative interpolation, and dynamic connectivity.

\chapter*{Engineering Selection Principles}
Use body-fitted grids for fixed geometry and high wall accuracy; immersed boundaries for complex moving geometry when exact wall stress is secondary; adaptive quadtrees or octrees for localized scales; moving meshes for moderate deformation; and overset grids for large relative rigid-body motion.

\chapter*{A Complete Computational Workflow}
A reliable simulation typically proceeds as follows:
\begin{enumerate}
\item choose the physical model from scale ratios such as $\mathrm{Kn}$, $\mathrm{Re}$, and $\mathrm{Ma}$;
\item define variables, constitutive laws, and boundary conditions;
\item choose a grid compatible with geometry and local scales;
\item select a conservative, stable, and sufficiently accurate discretization;
\item form the algebraic system through explicit or implicit advancement;
\item choose a solver and preconditioner from the matrix structure;
\item monitor model, grid, discretization, linear, and nonlinear errors separately.
\end{enumerate}

\chapter*{Conclusion}
\[
\boxed{\text{Liouville}\to\text{Vlasov / Boltzmann}\to\text{continuum equations}}
\]
\[
\boxed{\text{asymptotic limits and linearization}\to\text{numerical discretization}\to\text{stability, dissipation, and dispersion}}
\]
\[
\boxed{\text{sparse or dense matrices}\to\text{Krylov, LU, AMG, and MoM}\to\text{engineering grids}}
\]
The essence of computational mechanics is to decide which physical information should be retained, determine how that information should be discretized, and exploit operator and matrix structure to compute the discrete model efficiently and reliably.

\end{document}
